% titel: Blatt 3, A4: Ausbeute (Nitrobenzol aus Benzol)
% Teil II: der durchgerechnete Weg.

\subsection*{Gegeben und gesucht}
Reaktion: $\mathrm{C_6H_6} + \mathrm{HNO_3} \rightarrow
\mathrm{C_6H_5NO_2} + \mathrm{H_2O}$ (1:1). Gegeben: gewünscht
$m(\mathrm{C_6H_5NO_2}) = \qty{95}{g}$, reale Ausbeute
$\eta = \qty{76.3}{\percent}$. Gesucht: einzusetzende
Masse $m(\mathrm{C_6H_6})$ (Benzol).

Molare Massen (gerundet): $M(\mathrm{C_6H_6}) \approx \qty{78.11}{g/mol}$,
$M(\mathrm{C_6H_5NO_2}) \approx \qty{123.11}{g/mol}$.

\subsection*{Ansatz}
Reale Stoffmenge des Produkts bilden, durch $\eta$ teilen (theoretische
Stoffmenge), wegen 1:1 mit $M(\mathrm{C_6H_6})$ in die Benzol-Masse
umrechnen.

\subsection*{Rechnung}
\begin{align*}
  n_{\text{real}}(\mathrm{C_6H_5NO_2})
    &= \frac{95}{123.11} = 0.772\ \mathrm{mol}, \\
  n_{\text{theor}}(\mathrm{C_6H_5NO_2})
    &= \frac{0.772}{0.763} = 1.01\ \mathrm{mol}.
\end{align*}
Da das Verhältnis 1:1 ist, gilt
$n(\mathrm{C_6H_6}) = n_{\text{theor}} = 1.01\ \mathrm{mol}$:
\begin{equation}
  m(\mathrm{C_6H_6}) = 1.01 \times 78.11
  \approx \qty{79.0}{g} \approx \qty{79}{g}.
\end{equation}

\begin{ergebnis}
Es müssen $m(\mathrm{C_6H_6}) \approx \qty{79}{g}$ Benzol eingesetzt werden,
wenn die Ausbeute an Nitrobenzol $\qty{76.3}{\percent}$ beträgt.
\end{ergebnis}

\begin{kontrolle}
Theoretisch liefern $1.01\ \mathrm{mol}$ Benzol $1.01\ \mathrm{mol}$ Nitrobenzol
$= 124\ \mathrm{g}$; davon $76.3\%$ sind $\approx \qty{95}{g}$ --- stimmt.
Plausibilität: eingesetztes Benzol (79 g) ist deutlich weniger als die
theoretisch erhältliche Nitrobenzolmasse (124 g), weil Ausbeute und Molmassen
sich unterscheiden.
\end{kontrolle}

\begin{merksatz}
Mit einer unter 100\% liegenden Ausbeute gilt: real $= \eta\cdot$ theor, also
Eduktbedarf aus $\eta^{-1}$ rückwärts berechnen.
\end{merksatz}
